% Options for packages loaded elsewhere
\PassOptionsToPackage{unicode}{hyperref}
\PassOptionsToPackage{hyphens}{url}
\documentclass[
  10pt,
  letterpaper,
]{article}
\usepackage{xcolor}
\usepackage[margin=23mm]{geometry}
\usepackage{amsmath,amssymb}
\setcounter{secnumdepth}{5}
\usepackage{iftex}
\ifPDFTeX
  \usepackage[T1]{fontenc}
  \usepackage[utf8]{inputenc}
  \usepackage{textcomp} % provide euro and other symbols
\else % if luatex or xetex
  \usepackage{unicode-math} % this also loads fontspec
  \defaultfontfeatures{Scale=MatchLowercase}
  \defaultfontfeatures[\rmfamily]{Ligatures=TeX,Scale=1}
\fi
\usepackage{lmodern}
\ifPDFTeX\else
  % xetex/luatex font selection
\fi
% Use upquote if available, for straight quotes in verbatim environments
\IfFileExists{upquote.sty}{\usepackage{upquote}}{}
\IfFileExists{microtype.sty}{% use microtype if available
  \usepackage[]{microtype}
  \UseMicrotypeSet[protrusion]{basicmath} % disable protrusion for tt fonts
}{}
\makeatletter
\@ifundefined{KOMAClassName}{% if non-KOMA class
  \IfFileExists{parskip.sty}{%
    \usepackage{parskip}
  }{% else
    \setlength{\parindent}{0pt}
    \setlength{\parskip}{6pt plus 2pt minus 1pt}}
}{% if KOMA class
  \KOMAoptions{parskip=half}}
\makeatother
\usepackage{longtable,booktabs,array}
\newcounter{none} % for unnumbered tables
\usepackage{calc} % for calculating minipage widths
% Correct order of tables after \paragraph or \subparagraph
\usepackage{etoolbox}
\makeatletter
\patchcmd\longtable{\par}{\if@noskipsec\mbox{}\fi\par}{}{}
\makeatother
% Allow footnotes in longtable head/foot
\IfFileExists{footnotehyper.sty}{\usepackage{footnotehyper}}{\usepackage{footnote}}
\makesavenoteenv{longtable}
\setlength{\emergencystretch}{3em} % prevent overfull lines
\providecommand{\tightlist}{%
  \setlength{\itemsep}{0pt}\setlength{\parskip}{0pt}}
\usepackage[]{natbib}
\bibliographystyle{plainnat}
\usepackage{bookmark}
\IfFileExists{xurl.sty}{\usepackage{xurl}}{} % add URL line breaks if available
\urlstyle{same}
\hypersetup{
  pdftitle={Local Smoothing Does Not Establish Extrapolation: A Reproducible Toy Study},
  pdfauthor={Super Writer worked example},
  hidelinks,
  pdfcreator={LaTeX via pandoc}}

\title{Local Smoothing Does Not Establish Extrapolation: A Reproducible
Toy Study}
\author{Super Writer worked example}
\date{September 2026}

\begin{document}
\maketitle
\begin{abstract}
Better error on a familiar input range does not establish robustness
outside that range. We illustrate this distinction with an executable,
synthetic nearest-neighbor regression study. One-neighbor and
five-neighbor predictors are fitted to 40 noisy observations from a
known one-dimensional function, using five fixed training seeds. On an
interior evaluation grid, mean squared error is 0.1083 for one neighbor
and 0.0851 for five neighbors; on an extrapolation grid, the
corresponding errors are 2.6891 and 3.7012. These measurements support a
local comparison under the stated protocol, not a claim of universal
superiority or statistical significance. This is an educational writing
example, not a novel method or a submitted paper.
\end{abstract}

\section{Introduction}\label{introduction}

A results paragraph can be numerically correct and still promise more
than its experiment establishes. Consider a predictor that reduces error
on a fixed evaluation set. That observation answers a question about the
tested inputs, training procedure, and metric. It does not automatically
answer whether the predictor improves on a different domain, under a
different noise process, or with a different training budget. A useful
manuscript makes that distinction visible before interpreting the
result.

We use nearest-neighbor regression to make this writing problem
inspectable. The estimator is standard, so the reader need not
disentangle a new method from a new claim. Our question is deliberately
narrow: under one specified synthetic protocol, does averaging five
nearby training responses improve mean squared error both inside and
beyond the training input range?

The example has two deliverables. First, it provides a short experiment
that reconstructs every table entry from fixed seeds. Second, it shows
how a paper can report an improvement and an adverse result without
suppressing either. Neither deliverable is presented as a research
novelty or evidence that this writing workflow improves conference
acceptance rates.

\section{Background and Scope}\label{background-and-scope}

Nearest-neighbor regression predicts a response by averaging training
responses near a query. It is a conventional local estimator, discussed
alongside the bias-variance tradeoff in statistical learning texts
\citep{hastie2009}. We use this reference for the method's background,
not as support for the numerical results of this example. Those results
come only from the accompanying executable.

Increasing the neighborhood size changes the observations being
averaged. It can reduce sensitivity to individual noisy responses while
also smoothing local variation. The balance depends on the signal,
noise, sample size, and evaluation distribution. We therefore do not
begin with the assumption that five neighbors must be preferable. Both
neighborhood sizes and both evaluation domains are fixed in the script;
there is no automated hyperparameter search.

\section{Experimental Protocol}\label{experimental-protocol}

The underlying response is

\[
f(x) = 2x + \sin(4\pi x).
\]

For each seed in \(\{11,22,33,44,55\}\), we draw 40 training inputs
uniformly from \([0,1)\) and add independent uniform response noise in
\([-0.5,0.5]\). The random generator is Python's seeded
\texttt{random.Random}. Each training pair consumes one input draw
followed by one noise draw. The script specifies this order so
reproductions do not depend on an unstated sampling convention.

For \(k\in\{1,5\}\), prediction is the arithmetic mean of the responses
at the \(k\) closest training inputs, using absolute input distance.
Ties are resolved by training-row order. Both predictors see exactly the
same training set within a seed. We do not learn a distance metric,
rescale the input, or tune \(k\) using the evaluation results.

The in-domain grid is \(x_i=(i+0.5)/101\) for \(i=0,\ldots,100\). The
extrapolation grid is obtained by adding one to every grid point.
Evaluation targets are the noiseless function values \(f(x_i)\), rather
than new noisy observations. Thus the reported error measures recovery
of the chosen function, not prediction of future noisy labels.

For each training seed, method, and domain, we compute mean squared
error over the 101 grid points. We then report the arithmetic mean and
sample standard deviation across the five training seeds. The grid
points are not treated as 101 independent training replications. The
standard deviation describes variation across these five fits; it is not
a confidence interval.

\section{Results}\label{results}

{\def\LTcaptype{none} % do not increment counter
\begin{longtable}[]{@{}lrrr@{}}
\toprule\noalign{}
Evaluation domain & Neighbors & Mean MSE & Sample SD \\
\midrule\noalign{}
\endhead
\bottomrule\noalign{}
\endlastfoot
In-domain & 1 & 0.1083 & 0.0056 \\
In-domain & 5 & 0.0851 & 0.0128 \\
Extrapolation & 1 & 2.6891 & 1.4620 \\
Extrapolation & 5 & 3.7012 & 1.3386 \\
\end{longtable}
}

Inside the training input range, the five-neighbor predictor has lower
mean MSE than the one-neighbor predictor under this protocol. The
absolute reduction is approximately 0.0232. This supports the statement
that local averaging improves the observed mean error in this particular
comparison. We do not describe that reduction as statistically
significant: five fixed training replications and a descriptive table
are not an inferential test.

Beyond the training range, the ordering reverses. Five neighbors have
higher mean MSE, by approximately 1.0121. Omitting this domain would
leave the impression that the in-domain improvement establishes a
broader property than was tested. Reporting both domains instead makes
the claim's boundary explicit: the observed advantage does not carry
over to this extrapolation setting.

The per-seed CSV is part of the example, alongside the aggregation JSON
and the exact executable. Readers can inspect individual fits rather
than rely only on rounded table entries. Reproduction checks compare
unrounded values with a small floating-point tolerance and separately
check the displayed rounding; they do not infer scientific validity from
file existence.

\section{Interpretation}\label{interpretation}

For any query to the right of all training inputs, the ordering of input
distances is fixed. For a fixed training set, each nearest-neighbor
predictor therefore returns a constant throughout our extrapolation
interval: the response at the rightmost training input for \(k=1\), or
the mean response of the five rightmost inputs for \(k=5\). Neither
estimator extrapolates the chosen signal's linear trend or oscillations.
This follows from the estimator definition, independently of the
observed ranking in the results table.

That observation explains why extrapolation needs its own evaluation; it
does not prove that one neighbor will always outperform five there. The
ranking also depends on noisy boundary responses, the sampled training
inputs, and the target function. A different function or noise
realization could change it. The strongest defensible conclusion remains
a comparison of the reported errors under the stated conditions.

\section{Limitations}\label{limitations}

The study uses one function, one sample size, one noise distribution,
two neighborhood sizes, and five fixed seeds. It contains no real-world
data, external benchmark, validation-selected hyperparameters,
uncertainty-calibrated prediction intervals, or evidence about other
regressors. A smooth parametric baseline could answer a different and
useful question, but none was evaluated here. It must not appear in the
paper as a completed experiment.

The protocol was designed for an inspectable educational example, not
preregistered research. The example's author reviewed the resulting
table before writing this manuscript. We make no claim that the setup or
reported conditions were selected independently of the pedagogical
objective.

The manuscript was authored with AI assistance for this repository and
checked against its accompanying code and materials. Its original prose
and synthetic experiment are public teaching artifacts, not an
autonomous-writing benchmark, a blind human evaluation, or a paper
accepted by any venue.

\section{Conclusion}\label{conclusion}

The five-neighbor predictor improves mean error on the tested in-domain
grid but worsens it on the tested extrapolation grid. Keeping both
observations in the manuscript turns a broad claim of robustness into a
specific, reproducible comparison. The example illustrates an evidence
boundary; it does not introduce a new learning method or establish
general robustness.

\bibliography{references.bib}

\end{document}
